%% ====================================================================
\section{Experimental Design}
\label{sec:experiment}
%% ====================================================================

\subsection{Hypothesis}
\label{sec:hypothesis}

We formulate the investigation as a deterministic numerical verification with the
following decision framework.\footnote{Unlike a classical Neyman--Pearson test, there is
no randomness here---the computation is deterministic. We adopt the hypothesis language
for clarity of exposition, not to claim statistical inference.}

\paragraph{$H_0$ (Null hypothesis).}
The Riemann Hypothesis holds. All DQPTs in both System~1 and System~2 occur exclusively
at $\beta = 1/2$. Equivalently, for all $t \in \R$ and all $\beta \in (0,1)$ with
$\beta \neq 1/2$: $L(\beta, t) \neq 0$ and $G(\beta, t)$ does not exhibit a zero
corresponding to a zeta zero.

\paragraph{$H_1$ (Alternative hypothesis).}
The Riemann Hypothesis is false. There exists at least one DQPT at some
$\beta_0 \neq 1/2$, corresponding to a nontrivial zero of $\zeta(s)$ with
$\re(s) = \beta_0 \neq 1/2$.

We note that failure to reject $H_0$ does not constitute proof of~RH; it merely provides
additional numerical evidence consistent with~RH up to the scanned range.


\subsection{Test Protocol}
\label{sec:protocol}

The investigation proceeds in four stages.

\paragraph{Stage~1: Implementation Verification.}
Compute $L(1/2, t_k)$ and $G(1/2, t_k)$ at the first $100$ known nontrivial zeta zeros
($t_1 = 14.134725\ldots$, $t_2 = 21.022040\ldots$, etc.) and verify that they vanish
within numerical tolerance. This confirms that the implementation correctly reproduces
the theoretical predictions.

\paragraph{Stage~2: Critical-Line Zero Detection.}
Scan $t \in [0, T]$ at fixed $\beta = 1/2$, detecting zeros of $G(1/2, t)$ by sign
changes. Compare the detected zero positions with published
tables~\cite{odlyzko1989,lmfdb}. Verify that the zero count matches the Riemann--von
Mangoldt prediction $N(T)$.

\paragraph{Stage~3: Off-Critical-Line Scan (System~1 only).}
For each $\beta \in [0.1, 0.9]$ with spacing $\Delta\beta = 0.005$, excluding
$[0.4975, 0.5025]$, scan $t \in [0, T]$ and compute $|L(\beta, t)|$ using System~1 (the
eta function). System~2 (the $G$~function) is \emph{not} used for off-line scanning
because it can produce false positives at $\beta \neq 1/2$ due to phase-alignment zeros
(see Section~\ref{sec:system2}). Record the minimum value of $|L(\beta, t)|$ over the
$t$-range for each~$\beta$. If this minimum is below the zero-detection threshold~$\eps$,
flag it for further investigation (Stage~4). Apply the adaptive refinement strategy
described in Section~\ref{sec:params} (Grid resolution).

\paragraph{Stage~4: High-Precision Refinement.}
For any candidate $(\beta, t)$ flagged in Stage~3 with $|L(\beta, t)| < \eps$:
\begin{enumerate}
  \item[(a)] Increase $N$ and recompute to distinguish a genuine zero from a near-miss.
  \item[(b)] Refine $(\beta, t)$ by 2D Newton's method or simplex optimization.
  \item[(c)] If the refined $\beta$ is within tolerance of $1/2$, classify as a
             critical-line zero (consistent with~$H_0$). If $\beta$ remains bounded away
             from $1/2$, classify as a potential counterexample to~RH (rejecting~$H_0$).
\end{enumerate}


\subsection{Parameters}
\label{sec:params}

\paragraph{Truncation parameter $N$.}
\begin{itemize}
  \item \emph{Phase~1 (verification):} $N = 10^4$ with Euler--Maclaurin correction
        ($5$~terms). Justification: effective precision $\sim 10^{-38}$, far exceeding
        double-precision limits.
  \item \emph{Phase~3 (scanning):} $N = 10^3$ to $10^4$, depending on available compute.
        Justification: for a coarse scan, precision $\sim 10^{-2}$ suffices to detect
        zeros; flagged candidates are refined in Phase~4.
  \item \emph{Phase~4 (refinement):} $N = 10^5$ with Euler--Maclaurin correction.
\end{itemize}

\paragraph{Grid resolution.}
\begin{itemize}
  \item $\Delta\beta = 0.005$ for the coarse scan (at least $179$ off-critical-line
        $\beta$ values in $[0.1, 0.9]$ excluding a neighborhood of~$0.5$). The finer
        grid (compared to the naive $\Delta\beta = 0.05$) is necessary because a zero at
        $\beta_0 = 0.52$ would produce $|\eta(0.50 + it_0)| \sim 0.02$, which is too
        small to reliably detect at a grid point $0.02$ away. With $\Delta\beta = 0.005$,
        such a zero would be sampled at $\beta = 0.52$ (or within $0.0025$ of it), where
        $|\eta|$ is detectably small.
  \item $\Delta\beta = 0.001$ for refined scans near any flagged candidate.
  \item $\Delta t = 0.1$ for $t \in [0, 1000]$ ($10^4$ points per $\beta$ value).
  \item $\Delta t = 0.01$ for refined scans near flagged candidates.
\end{itemize}

\paragraph{Adaptive refinement strategy.}
Zeros with $|\beta_0 - 1/2| < \delta_{\mathrm{threshold}}$ may produce only marginal
dips in $|L(\beta, t)|$ at nearby grid points. To address this:
\begin{enumerate}
  \item \emph{Initial coarse scan:} Evaluate $|L(\beta, t)|$ on the
        $\Delta\beta = 0.005$ grid and record all local minima below a liberal threshold
        ($\eps_{\mathrm{coarse}} = 10^{-3}$).
  \item \emph{Adaptive refinement:} For each flagged $(\beta, t)$ region, bisect in both
        $\beta$ and~$t$ to find the true minimum of $|L|$. If the minimum decreases with
        refinement (indicating a genuine zero rather than a near-miss), continue refining
        until $|L| < 10^{-6}$ or the minimum stabilizes.
  \item \emph{Convergence check:} Increase $N$ and verify that $|L|$ continues to
        decrease (as expected for a genuine zero) rather than stabilizing (indicating a
        near-miss).
\end{enumerate}

\paragraph{Zero-detection tolerance.}
\begin{itemize}
  \item $\eps = 10^{-6}$ for classifying a point as a zero (coarse scan).
  \item $\eps = 10^{-12}$ for refined confirmation.
\end{itemize}

\paragraph{Scan range.}
\begin{itemize}
  \item $T_{\max} = 1000$ for the initial investigation (capturing $\sim 649$ zeros on
        the critical line).
  \item $T_{\max} = 10{,}000$ for extended runs ($\sim 10{,}142$ zeros).
  \item $T_{\max} = 100{,}000$ as an aspirational target with GPU acceleration
        ($\sim 138{,}069$ zeros).
\end{itemize}

\paragraph{Hardware.}
\begin{itemize}
  \item CPU: specified at runtime; any modern x86-64 processor suffices.
  \item GPU: NVIDIA RTX~2060 (1920 CUDA cores, 6~GB GDDR6 VRAM, Compute
        Capability~7.5).
  \item Precision: IEEE~754 double precision (64-bit, $\sim 15.9$ significant decimal
        digits).
\end{itemize}


\subsection{Validation}
\label{sec:validation}

\subsubsection{Unit Tests Against Reference Values}
\label{sec:unit-tests}

Each core mathematical function is tested against high-precision reference values
computed independently using \texttt{mpmath} (Python arbitrary-precision library) or
Mathematica. Specifically:
\begin{itemize}
  \item $S_N(s)$ for $N = 100$ and $s = 0.5 + 14.134725i$: compare to \texttt{mpmath}.
  \item $\vartheta(t)$ for $t = 14.134725$, $21.022040$, $100$, $1000$: compare to
        \texttt{mpmath}.
  \item $\eta(s)$ for $s = 0.5 + 14.134725i$: compare to \texttt{mpmath}.
  \item $Z_H(t)$ for $t = 14.134725$ (first zero): should be zero within precision.
\end{itemize}
Tolerance: agreement to $10$ significant digits (limited by double precision and~$N$).

\subsubsection{Cross-Validation: System~1 vs System~2}
\label{sec:cross-val}

Both systems must yield the same zero positions (to within numerical precision) at
$\beta = 1/2$. Specifically, if $L(1/2, t^*) = 0$, then $G(1/2, t^*)$ must also change
sign near $t = t^*$. We verify this for all detected zeros.

Tolerance: $|t^*_L - t^*_G| < 10^{-4}$ for each corresponding zero pair.

\subsubsection{Convergence Tests}
\label{sec:conv-tests}

For each of several fixed $(\beta, t)$ points, compute $L(\beta, t)$ for
$N = 10^2, 10^3, 10^4, 10^5$ and verify that the values converge. Specifically:
\begin{equation}\label{eq:conv-test}
  |L_N(\beta, t) - L_{N/10}(\beta, t)| < C \cdot (N/10)^{-\beta},
\end{equation}
where $C$ is a constant. The convergence rate should be consistent with the theoretical
$O(N^{-\beta})$.

\subsubsection{Comparison with Odlyzko's Tables}
\label{sec:odlyzko-compare}

The first $100$ detected zeros are compared with Odlyzko's published tables (available
from the LMFDB~\cite{lmfdb} and Andrew Odlyzko's
website~\cite{odlyzko1989,odlyzko2001}). Agreement is required to within the precision
of our computation.


\subsection{Controls}
\label{sec:controls}

The following consistency checks serve as controls to verify the correctness of the
implementation.

\paragraph{Control~1.}
At known zeros ($t = t_k$), verify:
\begin{equation}\label{eq:control1}
  F_1(1/2, t_k) = 1 \pm \eps,
\end{equation}
where $\eps = 0.05$ (accounting for finite-$N$ effects).\footnote{Under the Wei et~al.\
$\log$-base-$2$ normalization, this value would be $\ln 2 \approx 0.693$; we use the
natural log normalization throughout this document, giving
$F_1 = \beta_0 + 1 - \beta = 1$ at $\beta = \beta_0 = 1/2$.}

\paragraph{Control~2.}
Away from zeros, for generic~$t$ (e.g., $t = 50.0$, which is not near any zero), verify:
\begin{equation}\label{eq:control2}
  F_1(\beta, t) = (1 - \beta) \pm \eps
\end{equation}
for $\beta \in \{0.3, 0.5, 0.7\}$.\footnote{Under Wei et~al.\ normalization:
$(1 - \beta)\ln 2$.}

\paragraph{Control~3.}
Verify the zero count:
\begin{equation}\label{eq:control3}
  |N_{\mathrm{detected}}(T) - N_{\mathrm{predicted}}(T)| \le 1,
\end{equation}
where $N_{\mathrm{predicted}}(T) = \mathrm{round}((1/\pi)\vartheta(T) + 1)$. The
tolerance of~$1$ accounts for the $S(T) = (1/\pi)\arg\zeta(1/2 + iT)$ fluctuation term,
which can cause the actual zero count to differ from the smooth prediction by $\pm 1$ at
specific $T$ values.

\paragraph{Control~4.}
Verify that $G(1/2, t)$ is real-valued (imaginary part $< 10^{-14}$) for all sampled $t$
values, confirming the $Z$-function property.

\paragraph{Control~5.}
Verify the functional equation numerically: for selected $s$ values, confirm that
$\zeta(s)$ and $\zeta(1-s)$ satisfy the functional equation~\eqref{eq:func-eq} within
precision. (This tests the Riemann--Siegel theta function implementation indirectly.)
